\section{投资问题：价格在交易什么}

截至2026-07-22 13:52 CST，恒逸石化盘中价15.06元，日内涨幅约2.87\%，高／低15.43／14.33元，成交额约11.05亿元。过去约一年完整交易日的前复权区间为5.87--18.29元；当前较高点回撤约17.7\%，较低点上涨约156.6\%。核心判断是，价格已经从困境估值切换到景气兑现，当前更适合等待中报和三季报确认，而不是仅凭低静态PE追涨。

盘中成交量约7,371万股，五个完整交易日平均成交量约9,596万股／日、平均成交额约14.06亿元／日。成交活跃说明事件参与度高，却不能证明资金属性。本案没有形成可审计的龙虎榜机构席位、北向持仓变化、融资融券和基金净流入序列，因此不对“机构加仓”“北向看多”或“游资主导”作方向性结论。

\begin{exhibitbox}[二级市场状态]
\noindent\begin{tabularx}{\textwidth}{L{3.55cm}R{2.2cm}L{2.5cm}X}
\toprule
\textbf{指标} & \textbf{数值} & \textbf{状态} & \textbf{使用限制} \\
\midrule
盘中价／涨幅 & 15.06元／+2.87\% & 事件后高波动 & 非收盘价 \\
日内高／低 & 15.43／14.33元 & 波动较大 & 不能替代基本面区间 \\
盘中成交额 & 约11.05亿元 & 交易活跃 & 不识别资金主体 \\
五日VWAP & 14.6555元 & 市场成本锚 & 截至7月21日完整交易日 \\
一年QFQ高／低 & 18.29／5.87元 & 重估幅度大 & 前复权、历史不保证未来 \\
北向／机构／两融 & 未取得 & 权重为零 & 不用非官方值替代 \\
\bottomrule
\end{tabularx}
\end{exhibitbox}
\sourcenote{HY-MKT-QUOTE；HY-MKT-KLINE；data/verified\_market\_data.md，盘中截止2026-07-22 13:52，完整日线截止2026-07-21。}

\section{价格区间：技术观察与基本面价值分开}

第一支撑为13.0--13.5元，对应业绩预告后的验证区，也与基本面核心区间下沿相邻；第二支撑约12.0元。第一阻力16.3--17.1元，覆盖归一化Street锚、核心估值上沿与事件阻力；第二阻力约18.3元，为可观察年度高点。技术支撑与估值区间计算来源不同，不能因为数值接近就互相证明。

\begin{exhibitbox}[价格纪律与行动]
\centering
\begin{tikzpicture}[x=0.72cm,y=1cm]
\draw[very thick, steelgrey] (10.5,0) -- (22.5,0);
\draw[line width=5pt, riskred!35] (10.5,0) -- (13.0,0);
\draw[line width=5pt, riskgreen!45] (13.0,0) -- (13.5,0);
\draw[line width=5pt, accentblue!45] (13.5,0) -- (16.3,0);
\draw[line width=5pt, riskamber!55] (16.3,0) -- (16.5,0);
\draw[line width=5pt, riskred!45] (16.5,0) -- (22.5,0);
\foreach \x/\lab in {10.5/熊情景,12.0/第二支撑,13.5/观察下沿,15.06/当前价,15.7/目标价,16.5/重评阈值,18.3/年度高点,22.5/牛情景}
  {\draw[steelgrey] (\x,-0.12)--(\x,0.12); \node[font=\scriptsize, rotate=45, anchor=west] at (\x,-0.24) {\lab};}
\end{tikzpicture}
\end{exhibitbox}
\sourcenote{HY-MKT-KLINE；analysis/secondary\_market\_analysis.md；analysis/valuation\_model.md，截止2026-07-22。图中价格为研究阈值，不是自动交易指令。}

13.5元以下只有在裂解价差、现金回款和预测仍成立时才进入研究升级；跌破12元且基本面同步恶化，不作机械均值回归。13.5--16.3元维持观察；16.3--16.5元降低研究优先级；16.5元以上若没有利润和现金流上修，应触发重新评估。

\section{可转债事件：供给与降债的双面性}

剩余可转债本金约29.94亿元、转股价10.32元，当前股价明显高于转股价。公司此前阶段性决定不提前赎回，但2026-07-28后相关触发窗口可重新计算。若转股加速，市场短期面对股份供给，长期则可能减少债务和利息；不能只写“转股稀释”或“降债利好”。

二级市场应监控每日转股、季度转股公告、赎回决定和剩余本金。股本层面使用38.80694189亿股已披露稀释敏感性，避免在事件前低估分母；实际转股发生后，再按债务与财务费用更新。

\begin{keyinsight}[二级市场的“所以呢”]
价格已进入事件验证区，而不是低预期区。\textbf{行动上，等待中报现金流、三季报H2利润和7月28日后转债事件；不以成交放大推断机构态度，也不以均线替代基本面失效条件。}
\end{keyinsight}

\begin{riskbox}[交易风格失效条件]
若价格突破16.5元但盈利、现金与项目证据不变，预期透支加深；若跌破13.5元同时裂解、库存和现金恶化，则不再是单纯技术支撑。\textbf{行动上，价格和证据必须同向验证，不能只依据区间机械操作。}
\end{riskbox}
